\section{USAMO 2002}
        \begin{enumerate}
        	\item (\textit{USAMO 2002}) Let $S$ be a set with 2002 elements, and let $N$ be an integer with $0 \le N \le 2^{2002}$.  Prove that it is possible to color every subset of $S$ either blue or red so that the following conditions hold:


 (a) the union of any two red subsets is red;


 (b) the union of any two blue subsets is blue;


 (c) there are exactly $N$ red subsets.


 



\textbf{Solution 1}

Let a set colored in such a manner be \textit{properly colored}.  We prove that any set with $n$ elements can be properly colored for any $0 \le N \le 2^n$.  We proceed by induction.

 The base case, $n = 0$, is trivial.

 Suppose that our claim holds for $n = k$.  Let $s \in S$, $|S| = k + 1$, and let $S'$ denote the set of all elements of $S$ other than $s$.

 If $N \le 2^k$, then we may color all subsets of $S$ which contain $s$ blue, and we may properly color $S'$.  This is a proper coloring because the union of any two red sets must be a subset of $S'$, which is properly colored, and any the union of any two blue sets either must be in $S'$, which is properly colored, or must contain $s$ and therefore be blue.

 If $N > 2^k$, then we color all subsets containing $s$ red, and we color $N - 2^k$ elements of $S'$ red in such a way that $S'$ is colored properly.  Then $S$ is properly colored, using similar reasoning as before.  Thus the induction is complete.

 



\textbf{Solution 2}

If $N = 0$, color every subset blue.  If $N = 2^{2002}$, color every subset red.  Otherwise, let $S$ be $\{ 0, 1, \ldots , 2001 \}$.  Write $N$ in binary, i.e., let

 $N = \sum_{i=1}^{k} 2^{a_i}$,

 

 where each of the $a_i$ is an element of $S$.  We color each of the $a_i$ red and all the other elements of $S$ blue.  We color the empty set blue, and we color any other set the color of its largest element.  This satisfies the problem's first two conditions, as the largest element of the union of two red (or blue) sets will have a red (or blue) number as its largest element.  In addition, for each integer $n \in S$, there are $2^{n}$ subsets of $S$ with $n$ as a maximal element, so $\sum_{i=1}^{k} 2^{a_i} = N$ subsets of $S$ are colored red, as desired.

 
<i>Alternate solutions are always welcome. If you have a different, elegant solution to this problem, please \href{https://artofproblemsolving.com/wiki/index.php?title=2002\_USAMO\_Problems/Problem\_1&action=edit}{add it} to this page.</i>

 


	\item (\textit{USAMO 2002}) Let $ABC$ be a triangle such that


 
$\left( \cot \frac{A}{2} \right)^2 + \left( 2 \cot \frac{B}{2} \right)^2 + \left( 3 \cot \frac{C}{2} \right)^2 = \left( \frac{6s}{7r} \right)^2$,


 


 where $s$ and $r$ denote its \href{https://artofproblemsolving.com/wiki/index.php?title=Semiperimeter}{semiperimeter} and \href{https://artofproblemsolving.com/wiki/index.php?title=Inradius}{inradius}, respectively.  Prove that triangle $ABC$ is similar to a triangle $T$ whose side lengths are all positive integers with no common divisor and determine those integers.


 



\textbf{Solution}

Let $a,b,c$ denote $BC, CA, AB$, respectively.  Since the line from a triangle's incenter to one of its vertices bisects the angle at the triangle's vertex, the condition of the problem is equivalent to the

 $\left( \frac{s-a}{r} \right)^2 + 4\left( \frac{s-b}{r} \right)^2 + 9\left( \frac{s-c}{r} \right)^2 = \left( \frac{6s}{7r} \right)^2$,

 

 or

 $\frac{(s-a)^2}{36} + \frac{(s-b)^2}{9} + \frac{(s-c)^2}{4} = \frac{s^2}{36 + 9 + 4}$.

 

 But by the \href{https://artofproblemsolving.com/wiki/index.php?title=Cauchy-Schwarz\_Inequality}{Cauchy-Schwarz Inequality}, we know

 $\begin{matrix} (36 + 9 + 4) \left[ \frac{(s-a)^2}{36} + \frac{(s-b)^2}{9} + \frac{(s-c)^2}{4} \right] & \ge &\left[ (s-a) + (s-b) + (s-c) \right]^2\\ & = & s^2 \\ \qquad\qquad \quad \quad \frac{(s-a)^2}{36} + \frac{(s-b)^2}{9} + \frac{(s-c)^2}{4} & \ge &\frac{s^2}{36 + 9 + 4} \; , \end{matrix}$

 

 with equality only when $\frac{(s-a)^2}{36}, \frac{(s-b)^2}{9}, \frac{(s-c)^2}{4}$ are directly proportional to 36, 9, 4, respectively.  Therefore (clearing denominators and taking square roots) our problem requires that $(s-a), (s-b), (s-c)$ be directly proportional to 36, 9, 4, and since $a = (s-b) + (s-c)$ etc., this is equivalent to the condition that $a,b,c$ be in proportion with 13, 40, 45, Q.E.D.

 Sidenote: A 13, 40, 45 obtuse triangle has an integer area of 252 and an inradius of $\frac{36}{7}$.

 



\textbf{Solution 2 (Lagrange Multipliers)}

Let $D, E, F$ be the projections of the incenter onto the sides $BC$, $AC$, and $AB$, respectively, and let $x=AF . y=BD, z=CE$. We thus have

 
 \[\cot \frac{A}{2} = \frac{x}{r} , \cot \frac{B}{2} = \frac{y}{r} , \cot \frac{C}{2} = \frac{z}{r}\]

 and additionally, 

 
 \[x+y+z = s\]
 \[x^2 + 4 y^2 + 9 z^2 = (\frac{6s}{7})^2\]

 We first determine the side lengths of the triangle similar to this triangle with semiperimeter $1$. It is asked of us to show that such a triangle is unique, but we note that this is true if and only if the plane $x+y+z=1$ and the ellipsoid $x^2 + 4 y^2 + 9 z^2 = (\frac{6}{7})^2$ is tangent. This can be shown if we use the method of Lagrange Multipliers on the constraint $x+y+z=1$ and function $x^2 + 4 y^2 + 9 z^2$ to prove that $(\frac{6}{7})^2$ is a extrema.

 The Method of Lagrange Multipliers gives the system of 

 
 \[2x= \lambda\]
 \[8y= \lambda\]
 \[18z= \lambda\]
 \[x+y+z=1\]

 from which we easily obtain $(x,y,z)=(\frac{36}{49} , \frac{9}{49} , \frac{4}{49} )$. We calculate that this indeed satisfy $x^2 + 4 y^2 + 9 z^2 = (\frac{6}{7})^2$, so $(\frac{6}{7})^2$ is indeed an relative extrema. To check that it is a global extrema, one notes that as x, y, or z goes towards plus or minus infinity, the function $x^2 + 4 y^2 + 9 z^2$ goes to positive infinity, so that $(\frac{6}{7})^2$ is the global minima. Now, $AB=x+y=\frac{45}{49}$, $AC=x+z=\frac{40}{49}$, and $BC=y+z=\frac{13}{49}$, so we see that the desired side lengths are $AB=45, AC=40, BC=13$. $\square$. 

 
<i>Alternate solutions are always welcome. If you have a different, elegant solution to this problem, please \href{https://artofproblemsolving.com/wiki/index.php?title=2002\_USAMO\_Problems/Problem\_2&action=edit}{add it} to this page.</i>

 


	\item (\textit{USAMO 2002}) Prove that any monic polynomial (a polynomial with leading coefficient 1) of degree $n$ with real coefficients is the average of two monic polynomials of degree $n$ with $n$ real roots.


 



\textbf{Solution 1}

\textbf{Lemma.}  For any set of ordered pairs of reals $\{ (x_i , y_i ) \}_{i=1}^{n}$ with $x_i \neq x_j$ for all $i \neq j$, there exists a unique monic real polynomial  $P(x)$ of degree $n$ such that $P(x_i) = y_i$ for all integers $i \in [1,n]$.

 \textit{Proof 1.}  By the \href{https://artofproblemsolving.com/wiki/index.php?title=Lagrange\_Interpolation\_Formula}{Lagrange Interpolation Formula}, there exists a unique real polynomial $Q(x)$ of degree less than $n$ such that $Q(x_i) = y_i - x_i^n$.  Then it is necessary and sufficient that $P(x) = x^n + Q(x)$.

 \textit{Proof 2.}  Again by the Lagrange Interpolation Formula, there exists a unique real polynomial $Q(x)$ of degree less than $n$ such that $Q(x_i) = y_i$.  Then it is necessary and sufficient that $P(x) = Q(x) + \prod_{i=1}^{n}(x-x_i)$.

 Let $F(x)$ be the monic polynomial of degree $n$ given in the problem.  Let $x_1, \ldots, x_n$ be a sequence of strictly increasing reals, and choose $y_1, \ldots, y_n$ such that for odd $i$, $y_i > \max (0, 2F(x_i))$, but for even $i$, $y_i < \min (0, 2F(x_i))$.  Let $P(x)$ be the unique monic polynomial such that $P(x_i) = y_i$.  We shall prove that $P(x)$ and $Q(x) = 2F(x) - P(x)$ satisfy the problem's conditions.

 Note that for $i < n$, $P(x_i)$ and $P(x_{i+1})$ have different signs, so $P(x)$ has a real root between $x_i$ and $x_{i+1}$.  It follows that $P(x)$ has at least $n-1$ real roots.  But a polynomial with real coefficients must have an even number of non-real roots, so $P(x)$ must have $n$ real roots.  Similarly, $Q(x)$ must have $n$ real roots.  Q.E.D.

 



\textbf{Solution 2}

Let $F(x)$ be our polynomial.  If $n = 1$, then we may let $F = x + a$, which is the average of the polynomials $x$ and $x + 2a$, each of which has a real root.  Otherwise, let

 $G(x) = \prod_{i=1}^{n-1}(x-2i)$,

 

 and let

 $P(x) = x^n - kG(x)$

 

 and

 $Q(x) = 2F(x) - x^n + kG(x)$.

 

 We will prove that for sufficiently large $k$, $P(x)$ and $Q(x)$ satisfy the problem's conditions.

 We note that for the values $1, 3, \ldots, 2n-1$ of $x$, $G(x)$ alternates in sign, and always has magnitude at least 1 (since it is the product of several factors of magnitude at least one.  On the other hand, there exists some $c$ such that for all $x \in [1, 2n-1]$, $|x^n| < c$ and $|F(n) - x^n | < c$.  Let $k$ be greater than $c$.  Then both $P(x)$ and $Q(x)$ alternate in sign for $x = 1, 3, \ldots, 2n-1$.  It follows that each has at least $n-1$ real roots.  But since a polynomial with real coefficients can only have an even number of complex roots, this means that both $P$ and $Q$ must have $n$ roots.  Q.E.D.

 
<i>Alternate solutions are always welcome. If you have a different, elegant solution to this problem, please \href{https://artofproblemsolving.com/wiki/index.php?title=2002\_USAMO\_Problems/Problem\_3&action=edit}{add it} to this page.</i>

 


	\item (\textit{USAMO 2002}) Let $\mathbb{R}$ be the set of \href{https://artofproblemsolving.com/wiki/index.php?title=Real\_number}{real numbers}.  Determine all functions $f : \mathbb{R} \rightarrow \mathbb{R}$ such that


 
$f(x^2 - y^2) = xf(x) - yf(y)$


 


 for all pairs of real numbers $x$ and $y$.


 



\textbf{Solution 1}

We first prove that $f$ is \href{https://artofproblemsolving.com/wiki/index.php?title=Odd\_function}{ odd}.

 Note that $f(0) = f(x^2 - x^2) = xf(x) - xf(x) = 0$, and for nonzero $y$, $xf(x) + yf(-y) = f(x^2 - y^2) = xf(x) - yf(y)$, or $yf(-y) = -yf(y)$, which implies $f(-y) = -f(y)$.  Therefore $f$ is odd.  Henceforth, we shall assume that all variables are non-negative.

 If we let $y = 0$, then we obtain $f(x^2) = xf(x)$.  Therefore the problem's condition becomes

 $f(x^2 - y^2) + f(y^2) = f(x^2)$.

 

 But for any $a,b$, we may set $x = \sqrt{a}$, $y = \sqrt{b}$ to obtain

 $f(a-b) + f(b) = f(a)$.

 

 (It is well known that the only \href{https://artofproblemsolving.com/wiki/index.php?title=Continuous}{continuous} solutions to this functional equation are of the form $f(x) = kx$, but there do exist other solutions to this which are not solutions to the equation of this problem.)

 We may let $a = 2t$, $b = t$ to obtain $2f(t) = f(2t)$.

 Letting $x = t+1$ and $y = t$ in the original condition yields

 $\begin{matrix}f(2t+1) &=& (t+1)f(t+1) - tf(t) \qquad \\ &=& (t+1)[f(t) + f(1) ] - tf(t) \\ &=& f(t) + (t+1)f(1) \qquad \qquad \end{matrix}$

 

 But we know $f(2t + 1) = f(2t) + f(1) = 2f(t) + f(1)$, so we have $2f(t) + f(1) = f(t) + tf(1) + f(1)$, or

 $f(t) = tf(1)$.

 

 Hence all solutions to our equation are of the form $f(x) = kx$.  It is easy to see that real value of $k$ will suffice.

 



\textbf{Solution 2}

As in the first solution, we obtain the result that $f$ satisfies the condition

 $f(a) + f(b) = f(a+b)$.

 

 We note that

 $f(x) = f\left[ \left(\frac{x+1}{2}\right)^2 - \left( \frac{x-1}{2} \right)^2 \right] = \frac{x+1}{2} f \left( \frac{x+1}{2} \right) - \frac{x-1}{2} f \left( \frac{x-1}{2} \right)$.

 

 Since $f(2t) = 2f(t)$, this is equal to

 $\frac{(x+1)[f(x) +f(1)]}{4} - \frac{(x-1)[f(x) - f(1)]}{4} = \frac{xf(1) + f(x)}{2}$

 

 It follows that $f$ must be of the form $f(x) = kx$.

 


 <i>Alternate solutions are always welcome. If you have a different, elegant solution to this problem, please \href{https://artofproblemsolving.com/wiki/index.php?title=2002\_USAMO\_Problems/Problem\_4&action=edit}{add it} to this page.</i>

 


	\item (\textit{USAMO 2002}) Let $a, b$ be integers greater than 2.  Prove that there exists a positive integer $k$ and a finite sequence $n_1, n_2, \ldots, n_k$ of positive integers such that $n_1 = a$, $n_k = b$, and $n_in_{i+1}$ is divisible by $n_i + n_{i+1}$ for each $i$ ($1 \le i < k$).


 



\textbf{Solutions}

We may say that two integers $a$ and $b$ are \textit{connected} (and write $a \leftrightarrow b$) if there exists such a sequence of integers as described in the problem.  For reference, we note that $\leftrightarrow$ is an \href{https://artofproblemsolving.com/wiki/index.php?title=Equivalence\_relation}{equivalence relation}: it is \href{https://artofproblemsolving.com/wiki/index.php?title=Reflexive}{reflexive} ($a \leftrightarrow a$), \href{https://artofproblemsolving.com/wiki/index.php?title=Symmetric}{symmetric} ($a \leftrightarrow b$ implies $b \leftrightarrow a$), and \href{https://artofproblemsolving.com/wiki/index.php?title=Transitive}{transitive} ($a \leftrightarrow b \leftrightarrow c$ implies $a \leftrightarrow c$ ).

 



\textbf{Solution 1}

Note that for any divisor $d$ of some $n$, $n(d-1) + n = nd \mid n^2(d-1)$, so $n \leftrightarrow n(d-1)$.  It follows that in fact $n \leftrightarrow n(d-1)^k$, for any nonnegative integer $k$, and

 $n \leftrightarrow n(d-1) \leftrightarrow n(d-1)(d-2) \leftrightarrow \cdots \leftrightarrow n \prod_{i=c}^{d-1}i$

 

 for any positive integer $c < d$.

 For all integers $a > b > 2$, there exists some integer $\lambda$ such that $(b-1)^{\lambda} > a$.  Let

 $X = \prod_{i=b}^{(b-1)^{\lambda}}i$.

 

 We have

 $a \leftrightarrow \prod_{i=b}^{a}i \leftrightarrow (b-1)^{\lambda} \prod_{i=b}^{a}i \leftrightarrow (b-1)^{\lambda}  \prod_{i=b}^{a}i \cdot \prod_{i=a+1}^{(b-1)^{\lambda}-1} i = X$.

 

 But we also have

 $b \leftrightarrow b(b-1)^{\lambda} \leftrightarrow b (b-1)^{\lambda} \prod_{i=b+1}^{(b-1)^{\lambda}-1}i = X$.

 

 Hence $a \leftrightarrow X \leftrightarrow b$, so $a \leftrightarrow b$, as desired.

 



\textbf{Solution 2}

We note that for any integer $n \ge 3$, $n \leftrightarrow 2n$, for

 $n \leftrightarrow n(n-1) \leftrightarrow n(n-1)(n-2) \leftrightarrow n(n-2) \leftrightarrow 2n$

 

 It follows that for $n \ge 4$,

 $n \leftrightarrow n(n-1) \leftrightarrow n(n-1)(n-2) \leftrightarrow n(n-1)(n-2)(n-3) \leftrightarrow 2(n-1)(n-2) \leftrightarrow (n-1)(n-2) \leftrightarrow (n-1)$.

 

 Thus all integers greater than 2 are connected.

 



\textbf{Solution 3}

We note that $n_i + n_{i+1} \mid n_in_{i+1}$ if and only if

 $(n_i + n_{i+1}) \mid [(n_i + n_{i+1})n_i - n_in_{i+1}] = n_i^2$.

 

 Therefore for any divisor $d > n_i$ of $n_i^2$, $d - n_i \leftrightarrow n_i$.

 Now, for $a \ge 2$,

 $4a \leftrightarrow 4a^2 - 4a = 4a(a-1) \leftrightarrow 8a^2 - 4a(a-1) = 4a(a+1) \leftrightarrow 4(a+1)^2 - 4a(a+1) = 4(a+1)$.

 

 Also,

 $4 \leftrightarrow 4^2 - 4 = 12$.

 

 This means that all positive multiples of 4 are connected.

 Furthermore, for $a \ge 2$,

 $2a \leftrightarrow 2a(a-1)$,

 

 which implies that every even number greater than 2 is connected to a multiple of 4.

 Finally, for $a \ge 3$,

 $a \leftrightarrow a(a-1)$.

 

 Since $a(a-1)$ is even, this means that all integers greater than or equal to 2 are connected to some even number.

 Together, these imply that all integers greater than 2 are connected.

 



\textbf{Solution 4}

We note that if $d > 2$ is a divisor of $n$, then $n \leftrightarrow n(d-1)$.

 We say a positive integer $k$ is \textit{safe} if for all integers $n \ge 3$, $n \leftrightarrow kn$.  Note that the product of two safe numbers is also a safe number.  Define $f(n)$ ($n>2$) to be the smallest divisor of $n$ that is greater than 2.  We will prove that 2 is safe, by strong induction on $f(n)$.  For the case $f(n) = 3$, we have $n \leftrightarrow (3-1)n = 2n$.  If $f(n) > 3$, we note that $f[n(f(n) - 1)] < f(n)$, since $f(n) - 1$ must be less than all the divisors of $n$.  Thus the inductive hypothesis gives us $[f(n) - 1]n \leftrightarrow 2[f(n) - 1]n$.  Furthermore, we have $n \leftrightarrow [f(n)-1]n$ and $2n \leftrightarrow 2[f(n)-1]n$, both from our initial note.  Thus $n \leftrightarrow 2n$.

 We will now prove that every prime $p$ is safe, by strong induction.  We have already proven the base case $p = 2$.  Now, for odd $p$, $p+1$ is the product of odd primes less than $p$, so $p+1$ is safe.  Then we have

 $n \leftrightarrow (p+1)n \leftrightarrow p(p+1)n \leftrightarrow pn$.

 

 Thus the induction is complete, and all primes are safe.

 We have now shown that all integers greater than 1 are safe.  Specifically, $a$ and $b$ are safe, and $a \leftrightarrow ab \leftrightarrow b$.

 



\textbf{Solution 5}

Similarly to the first solution, we observe that for any integer $k < a$ and any positive integers $e_1, \ldots, e_k$,

 $a \leftrightarrow a \prod_{i=1}^{k}(a-i)^{e_i}$.

 

 We also note that if $a \leftrightarrow b$, then for any positive integer $k$, $ka \leftrightarrow kb$, for $(a+b) \mid (ab)$ implies $k(a+b) \mid k^2(ab)$.

 It is sufficient to prove that for any $a > 3$, $a \leftrightarrow (a-1)$.  If $a$ is composite, then there exist $m \ge n > 0$ such that $(a-1-m)(a-1-n) = a$, and

 $(a-1) \leftrightarrow (a-1-m)(a-1-n) \prod_{i=0}^{m-1}(a-1-i) = a \prod_{i=0}^{m-1}(a-1-i) \leftrightarrow a$.

 

 If, on the other hand, $a$ is prime, then we use strong induction.  For the base case, $a = 5$, we note $4 \leftrightarrow 3 \leftrightarrow 3(3-1) = 6 \leftrightarrow 5$.  Now, assuming that $a>5$ and that this holds for all primes less than $a$ (we know it holds for all composites), it is sufficient to show that $(a+1) \leftrightarrow (a-1)$, since $(a+1)$ is composite and therefore $(a+1) \leftrightarrow a$.  But since $(a+1)$ and $(a-1)$ are both even, it suffices to show that $\frac{a+1}{2} \leftrightarrow \frac{a-1}{2} = \frac{a+1}{2} - 1$.  But this is true, since $\frac{a+1}{2}$ either composite, or a prime less than $a$, and it is greater than 3, since $a > 5$.  Thus the induction is complete.

 


	\item (\textit{USAMO 2002}) 
I have an $n \times n$ sheet of stamps, from which I've been asked to tear out blocks of three adjacent stamps in a single row or column. (I can only tear along the perforations separating adjacent stamps, and each block must come out of the sheet in one piece.) Let $b(n)$ be the smallest number of blocks I can tear out and make it impossible to tear out any more blocks. Prove that there are real constants $c$ and $d$ such that 

 $\dfrac{1}{7} n^2 - cn \leq b(n) \leq \dfrac{1}{5} n^2 + dn$

  for all $n > 0$.


 



\textbf{Solution}

The upper bound requires an example of a set of $\frac{1}{5}n^2 + dn$ blocks whose removal makes it impossible to remove any further blocks. It suffices to show that we can tile the plane by tiles containing one block for every five stamps such that no more blocks can be chosen. Two such tilings are shown below with one tile outlined in heavy lines. Given an $n\times n$ section of the tiling take all blocks lying entirely within that section and add as many additional blocks as possible. If the basic tile is contained in an $(m + 1)\times (m + 1)$ square, then the $n\times n$ section is covered by tiles contained in a concentric $(n + 2m)\times (n + 2m)$ square. Hence there are at most $\frac{1}{5}(n + 2m)^2$ blocks entirely within the section. For an $n\times n$ section of the tiling, there are at most $4n$ blocks which lie partially in and partially out of that section (hence these blocks contain at most $8n$ stamps in the $n\times n$ square) and each of the additional blocks must contain one of these stamps. Thus there are at most $8n$ additional blocks. Thus there are at most
 \[\frac{1}{5}(n + 2m)^2 + 8n\leq \frac{1}{5}n^2 + \frac{4m^2 + 4m + 40}{5}n\]blocks total.

 
                    \begin{figure}[H]
                        \centering
                        \includegraphics[width=0.5\textwidth]{Solution/USAMO_I_2002_6_Solution_Figure_1.png}
                    \end{figure}
                    

 The lower bound requires an argument. Suppose that we have a set of $b(n)$ blocks whose removal makes removing any further blocks impossible. 

 



\textbf{Solution 1}

There are $2n(n - 2)$ potential blocks of three consecutive stamps in a row or column. Each of these must meet at least one of the $b(n)$ blocks removed. Conversely, each of the $b(n)$ blocks removed meets at most 14 of these potential blocks (5 oriented the same way, including itself, and 9 oriented the orthogonal way). Therefore $14b(n)\geq 2n(n - 2)$ or
 \[b(n)\geq\frac{1}{7}n^2 - \frac{2}{7}n.\]

 



\textbf{Solution 2}

Call a stamp used if it belongs to one of the $b(n)$ removed blocks. Consider the $(n - 2)^2$ five-stamp crosses centered at each stamp not on an edge of the sheet. Each cross must contain two used stamps. (One stamp not in the center is not enough to prevent another block from being torn out, and it is impossible to use one stamp in the center and use no other stamps in the cross.) In addition, each block not lying along an edge of the sheet lies entirely inside one cross, which thus contains three used stamps. There are at most $4n/3$ of the $b(n)$ blocks lying along the edges, hence there are at least $b(n) - 4n/3$ crosses containing three used stamps.

 Now count the number of pairs of a used stamp and a cross containing that stamp, in two ways. First counting block by block, we get $3b(n)$ used stamps, and each used stamp is contained in at most five crosses (exactly five if it is not on an edge), for a total of at most $15b(n)$ pairs. Next, counting cross by cross, each of the $(n - 2)^2$ crosses contains at least two used stamps and we have at least $b(n) - 4n/3$ crosses containing three used stamps, for a total of at least $2(n - 2)^2 + b(n) - 4n/3$ pairs. Therefore
 \[15b(n)\geq 2(n - 2)^2 + b(n) - \frac{4n}{3},\]or
 \[b(n)\geq\frac{1}{7}n^2 - \frac{16}{21}n.\]

 



\textbf{Solution 3}

Call a stamp used if it belongs to one of the $b(n)$ removed blocks. Count the number of pairs consisting of a used stamp and an adjacent unused stamp, in two ways.

 There are at least $(n - 2)^2 - 3b(n)$ unused stamps which are not on an edge. Since no more blocks can be torn out, either the stamp to the left or right and either the stamp or below such an unused stamp must be used. Thus we have at least $2n^2 - 8n - 6b(n)$ such pairs.

 Each block removed is adjacent to at most eight other stamps. However these eight stamps contain two blocks of three consecutive stamps. Hence at most six of these eight stamps can be unused. Thus each of the $b(n)$ blocks removed is involved in at most six pairs. Thus there are at most $6b(n)$ pairs.

 Combining these we have
 \[6b(n)\geq 2n^2 - 8n - 6b(n),\]or
 \[b(n)\geq\frac{1}{6}n^2 - \frac{2}{3}n.\]

 <i>Alternate solutions are always welcome. If you have a different, elegant solution to this problem, please \href{https://artofproblemsolving.com/wiki/index.php?title=2002\_USAMO\_Problems/Problem\_6&action=edit}{add it} to this page.</i>

 


\end{enumerate}